\documentclass[11pt,a4paper]{article}

\usepackage{fontspec}
\setmainfont{DejaVu Serif}[
  BoldFont={DejaVu Serif Bold},
  ItalicFont={DejaVu Serif},
  ItalicFeatures={FakeSlant=0.18},
  BoldItalicFont={DejaVu Serif Bold},
  BoldItalicFeatures={FakeSlant=0.18}
]
\setsansfont{DejaVu Sans}
\usepackage{polyglossia}
\setmainlanguage{english}
\usepackage[a4paper,margin=20.5mm]{geometry}
\usepackage{microtype}
\usepackage{setspace}
\setstretch{1.015}
\usepackage{amsmath,amssymb,amsthm,mathtools,bm}
\usepackage{booktabs,longtable,array,tabularx}
\usepackage{enumitem}
\usepackage{xcolor}
\usepackage{tikz}
\usetikzlibrary{arrows.meta,positioning}
\usepackage[unicode,hidelinks]{hyperref}
\usepackage[nameinlink,noabbrev]{cleveref}
\emergencystretch=3em
\allowdisplaybreaks

\hypersetup{
  pdftitle={Gear Contact Geometry as a Typed Finite Observation System},
  pdfauthor={Stas, Independent Research Program},
  pdfsubject={Finite contact cycles, frame quotients, contact dissipation, and rotation-number scope},
  pdfkeywords={gear contact, hunting tooth, finite torus, Willis relation, frame action, dissipation, rotation number, observation geometry}
}

\definecolor{obsblue}{RGB}{37,83,138}
\definecolor{obsred}{RGB}{150,45,45}
\definecolor{obsgreen}{RGB}{30,115,78}
\definecolor{lightblue}{RGB}{239,246,253}
\definecolor{lightred}{RGB}{253,242,242}
\definecolor{lightgreen}{RGB}{239,249,244}

\theoremstyle{definition}
\newtheorem{definition}{Definition}[section]
\newtheorem{model}[definition]{Model}
\theoremstyle{plain}
\newtheorem{theorem}[definition]{Theorem}
\newtheorem{proposition}[definition]{Proposition}
\newtheorem{corollary}[definition]{Corollary}
\theoremstyle{remark}
\newtheorem{remark}[definition]{Remark}
\newtheorem{warning}[definition]{Boundary}

\newcommand{\R}{\mathbb{R}}
\newcommand{\Z}{\mathbb{Z}}
\newcommand{\T}{\mathbb{T}}
\newcommand{\dd}{\mathop{}\!d}
\newcommand{\norm}[1]{\left\lVert #1\right\rVert}
\newcommand{\pair}[2]{\left\langle #1,#2\right\rangle}
\newcommand{\Ocal}{\mathcal O}
\newcommand{\Kcal}{\mathcal K}
\newcommand{\Rcal}{\mathcal R}
\newcommand{\Acal}{\mathcal A}
\newcommand{\Dcal}{\mathcal D}
\newcommand{\Pcal}{\mathcal P}
\newcommand{\one}{\bm 1}
\newcolumntype{L}[1]{>{\raggedright\arraybackslash}p{#1}}
\newcolumntype{Y}{>{\raggedright\arraybackslash}X}

\title{\bfseries Gear Contact Geometry\\
\large A typed finite observation system}
\author{Stas\\
\small Independent Research Program; ORCID: 0009-0000-2294-705X}
\date{Strict GO Core reference revision v1.1 --- 28 July 2026}

\begin{document}
\maketitle

\begin{center}
\fcolorbox{obsblue}{lightblue}{%
\parbox{0.92\linewidth}{\small
\textbf{Status.} This document supersedes the bilingual working draft
v1.0. It is migrated to GO Core v0.2 and the Frames--Forces--
Constraints--Dissipation Interface v0.1. The finite tooth-pair theorem
is exact for a declared pitch-event model. Contact loads, wear,
efficiency, lifetime, and tooth safety are not inferred from event
counts.}}
\end{center}

\begin{abstract}
Two ideal gears with positive tooth counts \(N_1,N_2\) admit an exact
finite event model on
\[
X_{N_1,N_2}=\Z_{N_1}\times\Z_{N_2},\qquad
T_\sigma(i,j)=(i+1,j+\sigma),
\]
where \(\sigma=-1\) denotes an external mesh and \(\sigma=+1\) an
internal mesh under the declared common angular orientation. Every
orbit has length \(\operatorname{lcm}(N_1,N_2)\), the state space has
\(\gcd(N_1,N_2)\) orbits, and one orbit covers every ordered tooth
pair exactly when the counts are coprime. We separate this exact
combinatorial layer from the active contact set, contact ratio,
normal load, sliding, dissipation, and material wear. The signed speed
ratio
\[
\frac{\omega_2}{\omega_1}
=\sigma\frac{N_1}{N_2}
\]
identifies only the coprime count ratio, not the common scale that
controls the cycle decomposition. For a planetary set, uniform
subtraction of frame angular speed is formalized as an
\((\R,+)\)-action; only a relative-observation protocol passes to its
quotient. For a periodic non-circular kinematic model, we define the
driver-period return map, its physical lift, invariant measures, and
the Poincare translation number. No KAM conclusion is drawn. The
result is a strict observation interface and a reproducible finite
benchmark, not a new gear-design or tribology theory.
\end{abstract}

\tableofcontents

\section{Scope and observation protocol}

The system descriptor is
\[
d=(N_1,N_2,\sigma,\mathfrak g,\mathfrak m,\vartheta),
\]
where \(N_a\in\mathbb N_{>0}\), \(\sigma\in\{-1,+1\}\),
\(\mathfrak g\) records the declared geometry and frame conventions,
\(\mathfrak m\) records the mesh class, and \(\vartheta\) contains any
constitutive parameters. The complete protocol is the typed chain
\[
X_t\xrightarrow{\Rcal_{\rm red}}\widetilde X_t
\xrightarrow{\Ocal_\lambda}Z_t
\xrightarrow{\Kcal_\lambda}\operatorname{Law}(Y_t\mid Z_t)
\xrightarrow{\widehat\theta_\lambda}\widehat\theta_t.
\]
The record \(\lambda\) must state the tooth labels, event clock,
reference frames, angular orientation, observation horizon, missed-
event and misclassification policy, load calibration, sampling rates,
noise model, and estimator.

\begin{warning}[Three layers that must not be identified]
The pitch-event state \(x_k\), the simultaneous active contact set
\(\Acal(t)\), and a material wear field are different objects. The
first is finite combinatorics. The second requires tooth geometry and
contact ratio. The third additionally requires loads, sliding,
material data, lubrication, and a constitutive evolution law.
\end{warning}

\section{Fixed-axis circular kinematics}

Let both angular coordinates be positive counterclockwise in one
declared inertial frame. Put
\[
\sigma=
\begin{cases}
-1,&\text{external mesh},\\
+1,&\text{internal mesh}.
\end{cases}
\]
For an ideal constant-ratio pair, the phase constraint is
\begin{equation}\label{eq:mesh-phase}
h(\theta_1,\theta_2)
=N_1\theta_1-\sigma N_2\theta_2-\phi_0
=0\pmod{2\pi}.
\end{equation}
On every smooth branch,
\begin{equation}\label{eq:mesh-velocity}
N_1\omega_1-\sigma N_2\omega_2=0,\qquad
\rho_{2\leftarrow1}
:=\frac{\omega_2}{\omega_1}
=\sigma\frac{N_1}{N_2},
\end{equation}
provided \(\omega_1\ne0\). The direction in
\(\rho_{2\leftarrow1}\) is part of the symbol; writing only ``gear
ratio'' is rejected.

\begin{proposition}[Ideal mesh-reaction power]\label{prop:ideal-power}
Let \(Q_c=\lambda\,\dd h\) be the reaction covector of the ideal
constraint \cref{eq:mesh-phase}. Then its total generalized power
vanishes on admissible velocities:
\[
\pair{Q_c}{(\omega_1,\omega_2)}
=\lambda(N_1\omega_1-\sigma N_2\omega_2)=0.
\]
\end{proposition}

This statement concerns the ideal constraint reaction. Sliding
friction, lubricant shear, impacts, backlash, bearing losses, and
profile errors belong to separate force and dissipation terms.

\section{Finite pitch-event dynamics}

\begin{definition}[Pitch-event state]
Choose an event section that records one reference tooth-pitch
crossing per advance of gear 1. The label state is
\[
X_{N_1,N_2}:=\Z_{N_1}\times\Z_{N_2}.
\]
Under \cref{eq:mesh-velocity}, its ideal event map is
\begin{equation}\label{eq:contact-shift}
T_\sigma(i,j)=(i+1,j+\sigma)
\quad\pmod{(N_1,N_2)}.
\end{equation}
\end{definition}

\begin{theorem}[Orbit length and cycle decomposition]
\label{thm:cycle-decomposition}
Every orbit of \(T_\sigma\) has length
\[
P=\operatorname{lcm}(N_1,N_2).
\]
The state space is the disjoint union of
\[
g=\gcd(N_1,N_2)
\]
directed cycles of length \(P\). A single orbit contains all
\(N_1N_2\) ordered tooth pairs if and only if \(g=1\).
\end{theorem}

\begin{proof}
The period is the order of \((1,\sigma)\) in
\(\Z_{N_1}\times\Z_{N_2}\). It is the least positive \(k\) satisfying
\[
k=0\pmod{N_1},\qquad \sigma k=0\pmod{N_2}.
\]
Because \(\sigma=\pm1\), this is
\(\operatorname{lcm}(N_1,N_2)\). Translation acts freely on each
orbit, so every orbit has this size. The orbit count is therefore
\[
\frac{N_1N_2}{\operatorname{lcm}(N_1,N_2)}
=\gcd(N_1,N_2).
\]
\end{proof}

\begin{corollary}[Hunting-tooth criterion]
Every tooth of gear 1 is paired with every tooth of gear 2 on one
event orbit exactly when \(\gcd(N_1,N_2)=1\).
\end{corollary}

The directed graph
\[
G_{N_1,N_2,\sigma}
=\bigl(X_{N_1,N_2},
\{(x,T_\sigma x):x\in X_{N_1,N_2}\}\bigr)
\]
is consequently a union of equal directed cycles. Label translations
\(\ell_{a,b}(i,j)=(i+a,j+b)\) commute with \(T_\sigma\), while
\(R(i,j)=(i,-j)\) conjugates \(T_\sigma\) to \(T_{-\sigma}\).
Thus \(P\), \(g\), and the cycle multiset are invariant under tooth
relabeling and orientation reversal. The ordered label sequence is
not invariant under those transformations.

\begin{warning}[Event-model boundary]
When transverse or overlap contact ratio exceeds one, several tooth
pairs may carry load simultaneously. The map \(T_\sigma\) then
records a chosen pitch event, not the full active contact set.
Undercutting, interference, flank modification, elastic deflection,
and loss of contact are also absent from \cref{eq:contact-shift}.
\end{warning}

\section{Ratio observation and identifiability}

Let
\[
\Ocal_{\rm ratio}(N_1,N_2,\sigma)
=\sigma\frac{N_1}{N_2}.
\]
This is a deterministic observation of the descriptor, not of an
individual tooth-pair state.

\begin{proposition}[Fiber of the magnitude-ratio observation]
\label{prop:ratio-fiber}
Suppose
\[
\left|\Ocal_{\rm ratio}\right|=\frac{u}{v},
\qquad \gcd(u,v)=1.
\]
Then every compatible positive tooth-count pair is uniquely
\[
(N_1,N_2)=(cu,cv),\qquad c\in\mathbb N_{>0}.
\]
For this pair,
\[
\gcd(N_1,N_2)=c,\qquad
\operatorname{lcm}(N_1,N_2)=cuv.
\]
Hence the speed-ratio magnitude does not identify contact-cycle count
or return period.
\end{proposition}

\begin{proof}
The equality \(N_1/N_2=u/v\) gives \(vN_1=uN_2\).
Coprimality implies \(u\mid N_1\) and \(v\mid N_2\), with one common
positive multiplier \(c\). The gcd and lcm formulas follow.
\end{proof}

On a bounded design class
\[
\Dcal_M=\{(N_1,N_2,\sigma):1\le N_1,N_2\le M\},
\]
the observation fiber is finite. On the unbounded class it is
countably infinite. Therefore an information-loss statement must
declare the descriptor domain; it cannot be inferred from a ratio
alone.

Useful deterministic readouts include
\begin{align*}
\Ocal_{\rm pair}(x_k)&=(i_k,j_k)
&&\text{only with calibrated tooth labels},\\
\Ocal_{\rm phase}(q)&=(\theta_1,\theta_2)\pmod{2\pi}
&&\text{with declared encoders and frames},\\
\Ocal_{\rm ratio}(d)&=\rho_{2\leftarrow1}
&&\text{with signed orientation},\\
\Ocal_{\rm active}(X_t)&=\Acal(t)
&&\text{with contact-resolution policy},\\
\Ocal_{\rm load}(X_t)&=(N_\alpha,T_\alpha)_{\alpha\in\Acal(t)}
&&\text{with force calibration}.
\end{align*}
None of these channels is silently substituted for another.

\section{Contact force, sliding, and dissipation}

For an active candidate pair \(\alpha=(i,j)\), let
\(g_\alpha(q)\ge0\) mean separation and let \(N_\alpha\ge0\) be the
compressive normal-force magnitude. In the smooth force regime,
\begin{equation}\label{eq:contact-complementarity}
g_\alpha\ge0,\qquad
N_\alpha\ge0,\qquad
N_\alpha g_\alpha=0.
\end{equation}
Let \(n_\alpha\) point from gear 1 toward gear 2, and let
\(F_\alpha=N_\alpha n_\alpha+T_\alpha\) be the force on gear 2. The
force on gear 1 is \(-F_\alpha\). With
\[
v_{{\rm rel},\alpha}=v_{2,\alpha}-v_{1,\alpha},\qquad
v_{T,\alpha}
=\bigl(I-n_\alpha n_\alpha^{\mathsf T}\bigr)v_{{\rm rel},\alpha},
\]
a dissipative tangential law satisfies
\begin{equation}\label{eq:friction-sign}
T_\alpha\cdot v_{T,\alpha}\le0,\qquad
P_{{\rm fric},\alpha}
:=-T_\alpha\cdot v_{T,\alpha}\ge0.
\end{equation}
If \(J_{{\rm rel},\alpha}\dot q=v_{{\rm rel},\alpha}\), the generalized
contact covector is
\begin{equation}\label{eq:force-pullback}
Q_\alpha=J_{{\rm rel},\alpha}^{\mathsf T}F_\alpha\in T_q^*Q.
\end{equation}
This pullback is required before contact forces are combined with
other generalized forces.

In a smooth active-set interval, the nonnegative total frictional loss
is
\begin{equation}\label{eq:friction-loss}
P_{\rm fric}
=-\sum_{\alpha\in\Acal(t)}
T_\alpha\cdot v_{T,\alpha}\ge0.
\end{equation}
For a stationary autonomous model with ideal normal constraints, its
energy balance has the interface form
\[
\dot E_L=P_{\rm act}-P_{\rm fric}-P_{\rm other},
\qquad P_{\rm other}\ge0.
\]
Moving supports, time-dependent tooth geometry, or imposed carrier
motion require the explicit support-power or
\(-\partial_tL\) terms of the common mechanics interface.

\begin{warning}[Nonsmooth boundary]
Contact onset, separation, impact, backlash switching, stick--slip
transitions, and inconsistent rigid Coulomb contact require a
nonsmooth inclusion or complementarity time-stepper. Existence,
uniqueness, and numerical convergence are not proved here.
\end{warning}

\section{Counts, exposure, dissipation, and wear}

For the pitch-event orbit \(x_k=T_\sigma^kx_0\), define the count
\[
C_{ij}(K)
=\sum_{k=0}^{K-1}\mathbf 1_{\{x_k=(i,j)\}}.
\]
For a resolved continuous active set, define instead
\begin{align}
I_{ij}(T)
&=\int_0^T
\mathbf 1_{\{(i,j)\in\Acal(t)\}}N_{ij}(t)\,\dd t,
\label{eq:normal-exposure}\\
D_{ij}(T)
&=\int_0^T
\mathbf 1_{\{(i,j)\in\Acal(t)\}}
\bigl[-T_{ij}(t)\cdot v_{T,ij}(t)\bigr]\,\dd t.
\label{eq:dissipative-exposure}
\end{align}
Here \(C_{ij}\) is a count, \(I_{ij}\) has impulse dimension, and
\(D_{ij}\) has energy dimension. They cannot be added or called
equivalent.

For any declared nonnegative channel \(w_\alpha\) with
\(W=\sum_\alpha w_\alpha>0\), put
\[
p_\alpha=\frac{w_\alpha}{W},\qquad
H_b(p)=-\sum_{\alpha:p_\alpha>0}p_\alpha\log_b p_\alpha,
\qquad b>1.
\]
If \(W=0\), the normalized distribution and its entropy are
undefined.

\begin{proposition}[Full-orbit event entropy]
For one full pitch-event orbit, the empirical distribution is uniform
on \(P=\operatorname{lcm}(N_1,N_2)\) states. Therefore
\[
H_b=\log_bP.
\]
Relative to the uniform distribution on all \(N_1N_2\) ordered pairs,
the reachability support defect is
\begin{equation}\label{eq:reachability-defect}
\Delta_b
=\log_b(N_1N_2)-\log_bP
=\log_b\gcd(N_1,N_2).
\end{equation}
For \(N_1N_2>1\), the normalized defect
\[
\overline\Delta
=\frac{\log_b\gcd(N_1,N_2)}{\log_b(N_1N_2)}
\]
is independent of \(b\).
\end{proposition}

Neither \(C_{ij}\), \(I_{ij}\), \(D_{ij}\), nor their entropies are a
wear law. A wear prediction additionally needs a declared
constitutive map with material hardness, local pressure, sliding
distance, lubrication, temperature, surface evolution, and
uncertainty.

\section{Fixed-axis gear trains}

Let a gear train have angular-velocity vector
\(\omega\in\R^m\). For every fixed-axis mesh edge \(e=(a,b)\), use
\[
N_a\omega_a-\sigma_eN_b\omega_b=0.
\]
For every rigid compound-shaft relation use
\(\omega_a-\omega_b=0\). Stacking these rows gives
\begin{equation}\label{eq:train-kernel}
C\omega=0.
\end{equation}
The admissible ideal velocity space is \(\ker C\); this formulation
does not assume in advance that it is one-dimensional.

Along an oriented simple mesh path
\(v_0\to v_1\to\cdots\to v_m\),
\begin{equation}\label{eq:path-ratio}
\frac{\omega_{v_m}}{\omega_{v_0}}
=\prod_{r=1}^{m}
\sigma_{e_r}\frac{N_{v_{r-1}}}{N_{v_r}},
\end{equation}
when all relevant velocities are nonzero. A closed ideal loop admits
a nonzero compatible motion only if the corresponding signed product
is one. Idler cancellation in a terminal ratio does not erase the
idler's tooth-contact history, bearing loads, inertia, or losses.

\section{Planetary frames as an explicit group action}

Let
\(\omega=(\omega_s,\omega_r,\omega_c)\in\R^3\) denote sun, ring, and
carrier angular velocities in one common axial orientation. A change
to a frame rotating at angular speed \(\alpha\) acts by
\begin{equation}\label{eq:frame-action}
G_\alpha(\omega)=\omega-\alpha\one,\qquad
\one=(1,1,1).
\end{equation}
It is an action of \((\R,+)\), because
\[
G_\alpha\circ G_\beta=G_{\alpha+\beta},\qquad
G_0=\operatorname{id},\qquad
G_\alpha^{-1}=G_{-\alpha}.
\]

\begin{theorem}[Relative quotient and Willis relation]
\label{thm:willis-quotient}
The map
\[
\pi:\R^3\longrightarrow\R^2,\qquad
\pi(\omega)
=(\omega_s-\omega_c,\omega_r-\omega_c)
\]
is constant exactly on the \(G_\alpha\)-orbits and induces
\[
\R^3/\operatorname{span}\{\one\}\cong\R^2.
\]
For a simple ideal planetary set, the Willis relation is the quotient
line
\begin{equation}\label{eq:willis}
N_s(\omega_s-\omega_c)
+N_r(\omega_r-\omega_c)=0,
\end{equation}
or, when the denominator is nonzero,
\[
\frac{\omega_s-\omega_c}{\omega_r-\omega_c}
=-\frac{N_r}{N_s}.
\]
\end{theorem}

\begin{proof}
\(\pi(\omega)=\pi(\omega')\) exactly when
\(\omega'-\omega=\alpha\one\) for some \(\alpha\). Thus its fibers are
the action orbits and its kernel is
\(\operatorname{span}\{\one\}\). Equation \cref{eq:willis} depends
only on \(\pi(\omega)\).
\end{proof}

Uniform frame subtraction is therefore a physical change of
representation. It becomes an observational gauge redundancy only
for a protocol of the form
\(\Ocal_{\rm rel}=f\circ\pi\), which deliberately discards common
absolute spin. A laboratory encoder that measures \(\omega_c\)
separately does not pass to this quotient.

If the ring is fixed in the laboratory, \(\omega_r=0\), then
\[
\omega_c=\frac{N_s}{N_s+N_r}\,\omega_s.
\]
For \(N_s=30\), \(N_r=78\), and
\(\omega_s=1\ {\rm rad\,s^{-1}}\), this gives
\(\omega_c=5/18\ {\rm rad\,s^{-1}}\). For a standard equal-module
sun--planet--ring geometry, the separate compatibility condition is
\(N_r=N_s+2N_p\), giving \(N_p=24\) in this example.

\section{Periodic non-circular kinematics}

Let \(r_1,r_2:S^1\to(0,\infty)\) be positive \(C^1\) effective pitch
radii, and assume the driver angle \(\theta_1\) is monotone. A
kinematic rolling model is
\begin{equation}\label{eq:noncircular-ode}
\frac{\dd\theta_2}{\dd\theta_1}
=\sigma\frac{r_1(\theta_1)}{r_2(\theta_2)}.
\end{equation}
This equation is a model input. It does not prove manufacturability,
absence of interference, force closure, or elastic stability.

For the lifted solution
\(\widetilde\theta_2(\theta_1;x)\) with initial value \(x\), define
the one-driver-turn return lift
\begin{equation}\label{eq:return-lift}
\widetilde F(x)
=\widetilde\theta_2(2\pi;x),\qquad
\widetilde F(x+2\pi)=\widetilde F(x)+2\pi.
\end{equation}
Uniqueness of the scalar \(C^1\) flow makes the induced
\(F:S^1\to S^1\) orientation-preserving, even when the output angle
decreases along the driven orbit.

\begin{theorem}[Translation number with declared lift]
\label{thm:translation-number}
For the orientation-preserving circle homeomorphism \(F\) and the
physical lift \(\widetilde F\), the limit
\begin{equation}\label{eq:translation-number}
\tau(\widetilde F)
=\lim_{n\to\infty}
\frac{\widetilde F^n(x)-x}{2\pi n}
\end{equation}
exists and is independent of \(x\). Replacing the lift by
\(\widetilde F+2\pi k\) changes \(\tau\) by \(k\); hence
\(\rho(F)=\tau(\widetilde F)\bmod\Z\) is lift-independent.
If \(\mu\) is any \(F\)-invariant Borel probability measure, then
\[
\tau(\widetilde F)
=\frac1{2\pi}\int_{S^1}
\bigl(\widetilde F(\widetilde x)-\widetilde x\bigr)\,\dd\mu(x),
\]
where the displacement is the well-defined periodic function induced
by the selected lift.
\end{theorem}

For constant radii,
\(\widetilde F(x)=x+2\pi\sigma r_1/r_2\), so
\(\tau=\sigma r_1/r_2\). Elasticity, backlash, impact, or controller
hysteresis may destroy the single-valued homeomorphism hypothesis.
No KAM persistence, Diophantine stability, or locking theorem follows
from \cref{eq:noncircular-ode} alone.

\section{Reproducible benchmarks and machine gates}

\begin{center}
\begin{tabular}{rrrrr}
\toprule
\((N_1,N_2)\) & \(\gcd\) & \(\operatorname{lcm}\) &
orbit fraction & \(\lvert\rho_{2\leftarrow1}\rvert\)\\
\midrule
\((12,25)\) & \(1\) & \(300\) & \(1\) & \(12/25\)\\
\((12,18)\) & \(6\) & \(36\) & \(1/6\) & \(2/3\)\\
\((2,3)\) & \(1\) & \(6\) & \(1\) & \(2/3\)\\
\((4,6)\) & \(2\) & \(12\) & \(1/2\) & \(2/3\)\\
\bottomrule
\end{tabular}
\end{center}

The release tests verify the orbit theorem by exhaustive enumeration
over a finite tooth-count range, ratio fibers, graph conjugacy,
explicit entropy bases, contact complementarity, dissipation signs,
force-power covariance, train path products, the planetary group
action and quotient, the Willis example, and translation-number
convergence for benchmark circle maps.

\begin{longtable}{L{33mm}L{48mm}L{62mm}}
\toprule
Gate & Required declaration & Rejected inference\\
\midrule
\endhead
Event semantics & pitch section and label convention &
event pair equals simultaneous contact set\\
Ratio direction & \(\omega_2/\omega_1\), sign, and frame &
unnamed ``gear ratio''\\
Frame action & \(G_\alpha\) and observed quotient &
calling every carrier subtraction a gauge\\
Contact force & gap sign, active set, force-on-body convention &
unsigned or double-counted reactions\\
Dissipation & \(P_{\rm diss}\ge0\) with relative velocity &
event count used as energy loss\\
Exposure & count, impulse, and energy kept dimensionally distinct &
load count called wear\\
Entropy & base, support, positive-total guard &
\(\log\) without base or zero normalization\\
Rotation number & return map, lift, and hypotheses &
KAM language from a local ratio equation\\
\bottomrule
\end{longtable}

\section{Claim boundary}

\begin{center}
\begin{tabularx}{\linewidth}{L{42mm}L{28mm}Y}
\toprule
Statement & Status & Required hypotheses\\
\midrule
Finite orbit decomposition & theorem &
positive integer tooth counts; pitch-event map
\cref{eq:contact-shift}\\
Ratio fiber & proposition &
exact constant ratio; positive tooth counts\\
Ideal reaction power & proposition &
smooth ideal holonomic constraint\\
Willis quotient & theorem &
common axis convention; ideal simple planetary kinematics\\
Translation number & classical theorem &
orientation-preserving circle homeomorphism and selected lift\\
Contact friction & constitutive model &
declared active set, force convention, smooth force regime\\
Wear prediction & absent &
requires an external validated material-evolution model\\
\bottomrule
\end{tabularx}
\end{center}

This revision establishes no new gear-design standard, load-capacity
formula, lubrication model, lifetime law, or safety result. Its exact
mathematics is elementary finite-group dynamics plus classical circle
rotation theory. The contribution is a strict separation of
descriptor, event, contact, frame, dissipation, and observation
layers suitable for GO Core validation.

\begingroup
\small
\hbadness=10000
\begin{thebibliography}{9}

\bibitem{iso21771}
International Organization for Standardization,
\emph{ISO 21771-1:2024, Cylindrical involute gears and gear pairs --
Part 1: Concepts and geometry}, 2024.
\href{https://www.iso.org/standard/84949.html}
{ISO 21771-1:2024}.

\bibitem{litvin}
F.~L. Litvin and A. Fuentes,
\emph{Gear Geometry and Applied Theory}, 2nd ed.,
Cambridge University Press, 2004.
\href{https://doi.org/10.1017/CBO9780511547126}
{doi:10.1017/CBO9780511547126}.

\bibitem{noncircular}
F.~L. Litvin, A. Fuentes-Aznar, I. Gonzalez-Perez, and K. Hayasaka,
\emph{Noncircular Gears: Design and Generation},
Cambridge University Press, 2009, ISBN 978-0-521-76170-3.

\bibitem{arnold}
V.~I. Arnold, \emph{Mathematical Methods of Classical Mechanics},
2nd ed., Springer, 1989.
\href{https://doi.org/10.1007/978-1-4757-2063-1}
{doi:10.1007/978-1-4757-2063-1}.

\bibitem{marsden}
J.~E. Marsden and T.~S. Ratiu,
\emph{Introduction to Mechanics and Symmetry}, 2nd ed.,
Springer, 1999.
\href{https://doi.org/10.1007/978-0-387-21792-5}
{doi:10.1007/978-0-387-21792-5}.

\end{thebibliography}
\endgroup

\end{document}
